\documentclass[aspectratio=169, 10pt]{beamer}

% ==============================================================================
% THEME & COLOR CONFIGURATION (Professional Academic Navy-Teal Palette)
% ==============================================================================
\usetheme{Madrid}
\usefonttheme{professionalfonts}

\definecolor{NavyPrimary}{RGB}{15, 44, 89}      % #0F2C59
\definecolor{TealSecondary}{RGB}{2, 132, 199}   % #0284C7
\definecolor{DarkSlate}{RGB}{30, 41, 59}        % #1E293B
\definecolor{LightBg}{RGB}{248, 250, 252}       % #F8FAFC
\definecolor{AccentRed}{RGB}{220, 38, 38}       % #DC2626
\definecolor{AccentGreen}{RGB}{22, 163, 74}     % #16A34A
\definecolor{CardBg}{RGB}{241, 245, 249}        % #F1F5F9

\setbeamercolor{palette primary}{bg=NavyPrimary, fg=white}
\setbeamercolor{palette secondary}{bg=TealSecondary, fg=white}
\setbeamercolor{palette tertiary}{bg=DarkSlate, fg=white}
\setbeamercolor{structure}{fg=NavyPrimary}
\setbeamercolor{titlelike}{parent=palette primary}
\setbeamercolor{block title}{bg=NavyPrimary, fg=white}
\setbeamercolor{block body}{bg=CardBg, fg=DarkSlate}
\setbeamercolor{block title alerted}{bg=AccentRed, fg=white}
\setbeamercolor{block body alerted}{bg=LightBg, fg=DarkSlate}
\setbeamercolor{block title example}{bg=AccentGreen, fg=white}
\setbeamercolor{block body example}{bg=LightBg, fg=DarkSlate}

% ==============================================================================
% REQUIRED PACKAGES
% ==============================================================================
\usepackage[utf8]{inputenc}
\usepackage{amsmath, amssymb, amsfonts}
\usepackage{booktabs}
\usepackage{array}
\usepackage{multicol}
\usepackage{tikz}
\usepackage{microtype}

% Customize header/footer
\setbeamertemplate{navigation symbols}{}
\setbeamertemplate{footline}{
  \leavevmode%
  \hbox{%
  \begin{beamercolorbox}[wd=.333333\paperwidth,ht=2.25ex,dp=1ex,center]{author in head/foot}%
    \usebeamerfont{author in head/foot}Lithium-Ion Battery RUL Framework
  \end{beamercolorbox}%
  \begin{beamercolorbox}[wd=.333333\paperwidth,ht=2.25ex,dp=1ex,center]{title in head/foot}%
    \usebeamerfont{title in head/foot}Project Review 1: Version 1.0
  \end{beamercolorbox}%
  \begin{beamercolorbox}[wd=.333333\paperwidth,ht=2.25ex,dp=1ex,right]{date in head/foot}%
    \usebeamerfont{date in head/foot}\insertframenumber{} / \inserttotalframenumber\hspace*{2ex} 
  \end{beamercolorbox}}%
  \vskip0pt%
}

% ==============================================================================
% PRESENTATION METADATA
% ==============================================================================
\title[Battery RUL Prediction Framework]{Intelligent Lithium-Ion Battery Remaining Useful Life (RUL) Prediction Framework}
\subtitle{Machine Learning Pipeline Benchmarking, 5-Fold Cross-Validation, and Hyperparameter Optimization}
\author[Project Team]{Academic Project Review 1}
\institute[Engineering Department]{Department of Machine Learning \& Data Science}
\date[\today]{Semester Project Presentation | Version 1.0}

% ==============================================================================
% DOCUMENT BEGINS
% ==============================================================================
\begin{document}

% ------------------------------------------------------------------------------
% SLIDE 1: TITLE SLIDE
% ------------------------------------------------------------------------------
\begin{frame}[plain]
  \titlepage
  \begin{center}
    \footnotesize
    \textcolor{DarkSlate}{\textbf{Focus:} Data-Driven Modeling $\bullet$ GPU Acceleration $\bullet$ 10-Model Empirical Benchmark $\bullet$ Diagnostic Analytics}
  \end{center}
\end{frame}

% ------------------------------------------------------------------------------
% SLIDE 2: EXECUTIVE SUMMARY & PROJECT PHASING
% ------------------------------------------------------------------------------
\begin{frame}{Executive Summary \& Project Phasing}
\begin{columns}[T]
  \begin{column}{0.48\textwidth}
    \begin{block}{\textbf{Version 1.0: Current Implementation (Review 1)}}
      \begin{itemize}\small
        \item \textbf{Dataset Processing}: 14,992 cycles analyzed across 7 operational degradation features (0 nulls, 0 duplicates).
        \item \textbf{10-Model ML Benchmark}: Evaluated bagging, boosting, instance-based, trees, and linear baselines.
        \item \textbf{GPU Acceleration}: CUDA-accelerated XGBoost and multi-core CPU parallelization.
        \item \textbf{Rigorous 5-Fold CV}: Systematic hyperparameter tuning via \texttt{GridSearchCV}.
        \item \textbf{Top Result}: Extra Trees Regressor achieved \textbf{$R^2 = 0.9971$} and \textbf{MAE = 7.71 cycles}.
      \end{itemize}
    \end{block}
  \end{column}
  
  \begin{column}{0.48\textwidth}
    \begin{block}{\textbf{Version 2.0: Future Roadmap (Review 2)}}
      \begin{itemize}\small
        \item \textbf{Battery Digital Twin}: Integration of optimized ML models with live BMS telemetry streams.
        \item \textbf{Real-Time State Estimation}: Continuous State of Health (SOH) and RUL tracking with dynamic confidence intervals.
        \item \textbf{Online Incremental Learning}: Recursive weight updates under variable thermal and discharge stress.
        \item \textbf{Closed-Loop Control}: Automated charging rate optimization to maximize cell longevity.
      \end{itemize}
    \end{block}
  \end{column}
\end{columns}
\end{frame}

% ------------------------------------------------------------------------------
% SLIDE 3: PROBLEM STATEMENT & BATTERY ELECTROCHEMISTRY
% ------------------------------------------------------------------------------
\begin{frame}{Problem Statement: Non-Linear Battery Aging}
\begin{columns}[T]
  \begin{column}{0.55\textwidth}
    \textbf{Electrochemical Degradation Mechanisms:}
    \begin{itemize}\small
      \item \textbf{Loss of Lithium Inventory (LLI)}: Irreversible electrolyte decomposition continuously thickens the Solid Electrolyte Interphase (SEI) film on the anode.
      \item \textbf{Loss of Active Material (LAM)}: Cyclic intercalation stress induces mechanical micro-cracking and particle detachment at the cathode.
      \item \textbf{Impedance Rise ($R_{\text{int}}$)}: Polarization escalation triggers terminal voltage drops during discharge and premature charge cutoffs.
    \end{itemize}
    \vspace{0.5em}
    \begin{alertblock}{\textbf{End-of-Life (EOL) Criterion}}
      \small When battery capacity drops below \textbf{70\%--80\% of nominal rating}, the cell enters a steep degradation knee prone to thermal runaway.
    \end{alertblock}
  \end{column}

  \begin{column}{0.42\textwidth}
    \begin{block}{\textbf{Mathematical Formulation}}
      \small
      Given operational cycle index $t$, the Remaining Useful Life $\text{RUL}(t)$ is defined as:
      $$\text{RUL}(t) = N_{\text{EOL}} - N(t)$$
      where $N_{\text{EOL}}$ denotes the cycle count where capacity crosses the failure threshold.
    \end{block}
    \vspace{0.3em}
    \textbf{Core Engineering Challenge:}
    \begin{itemize}\footnotesize
      \item Estimate $\text{RUL}(t)$ accurately and non-invasively using external surface sensor telemetry (voltage cutoffs, discharge duration, dwell times).
    \end{itemize}
  \end{column}
\end{columns}
\end{frame}

% ------------------------------------------------------------------------------
% SLIDE 4: PROJECT OBJECTIVES (VERSION 1.0 SCOPE)
% ------------------------------------------------------------------------------
\begin{frame}{Project Objectives (Version 1.0 Scope)}
\begin{columns}[T]
  \begin{column}{0.5\textwidth}
    \begin{block}{\textbf{Objective 1: Statistical EDA \& Integrity}}
      \small Fully evaluate statistical moments, probability distributions, skewness, kurtosis, and correlation structure across 14,992 degradation records.
    \end{block}
    \vspace{0.3em}
    \begin{block}{\textbf{Objective 2: Standardized Pipeline Design}}
      \small Build leak-free Scikit-Learn pipelines encapsulating \texttt{StandardScaler} with regression estimators to eliminate data leakage.
    \end{block}
    \vspace{0.3em}
    \begin{block}{\textbf{Objective 3: Hardware Acceleration}}
      \small Deploy NVIDIA CUDA GPU computing for high-throughput tree boosting (XGBoost GPU) and multi-core CPU parallel execution.
    \end{block}
  \end{column}

  \begin{column}{0.48\textwidth}
    \begin{block}{\textbf{Objective 4: Systematic 5-Fold Tuning}}
      \small Perform exhaustive parameter sweeps across 10 diverse ML models using 5-fold cross-validation (\texttt{GridSearchCV}) scoring for negative MAE.
    \end{block}
    \vspace{0.3em}
    \begin{block}{\textbf{Objective 5: Multi-Model Benchmarking}}
      \small Quantitatively rank all 10 baseline and tuned models across Train/Test MAE, RMSE, $R^2$, and inference execution latencies.
    \end{block}
    \vspace{0.3em}
    \begin{block}{\textbf{Objective 6: Diagnostic Visual Analytics}}
      \small Develop high-precision parity bounds, residual homoscedasticity, error CDF curves, and feature importance rankings.
    \end{block}
  \end{column}
\end{columns}
\end{frame}

% ------------------------------------------------------------------------------
% SLIDE 5: COMPREHENSIVE LITERATURE REVIEW
% ------------------------------------------------------------------------------
\begin{frame}{Literature Review: 5 Foundational Research Papers}
\small
\begin{table}[]
\centering
\resizebox{\textwidth}{!}{%
\begin{tabular}{@{}llll@{}}
\toprule
\textbf{Paper / Reference} & \textbf{Publication Venue} & \textbf{Key Methodology / Focus} & \textbf{Takeaway for Our Project} \\ \midrule
\textbf{1. Haghighi et al. (2026)} & Elsevier \textit{Results in Eng.} & Stacked LSTM, TCN, Conformer Ensembles & Validates multi-model evaluation and ensemble stability. \\
\textbf{2. Wu et al. (2026)} & MDPI \textit{Sensors} & Integrated Ensembles (RF, RVM, Elastic Net) & Proves tree ensembles achieve high accuracy with fast inference. \\
\textbf{3. Liu \& Chen (2019)} & \textit{IEEE Access} & Health Indicators + Gaussian Process Reg. & Validates indirect voltage/charging indicators as RUL proxies. \\
\textbf{4. Wang et al. (2023)} & \textit{IEEE / arXiv} & Two-Stage Early Non-Linear Prognostics & Confirms non-linear tree models handle steep degradation knees. \\
\textbf{5. Zhang et al. (2024)} & \textit{IEEE / arXiv} & Hybrid Attention Networks vs. Tabular Noise & Shows tuned tree ensembles excel on tabular cycle features. \\ \bottomrule
\end{tabular}%
}
\end{table}
\vspace{0.5em}
\begin{block}{\textbf{Synthesis of Literature Findings}}
  \footnotesize
  \textbf{Key Finding:} On tabular cycle metrics, tuned tree ensembles (Extra Trees, Random Forest, XGBoost, LightGBM) achieve competitive or superior accuracy compared to deep neural networks while maintaining sub-millisecond execution latencies suitable for edge controllers.
\end{block}
\end{frame}

% ------------------------------------------------------------------------------
% SLIDE 6: RESEARCH GAP IDENTIFICATION
% ------------------------------------------------------------------------------
\begin{frame}{Research Gap Analysis}
\begin{columns}[T]
  \begin{column}{0.48\textwidth}
    \begin{block}{\textbf{Gap 1: Computational Overhead on Edge BMS}}
      \small Deep learning models (Transformers, LSTMs) demand high parameter overhead and GPU memory. A rigorous benchmark of lightweight, high-accuracy tree algorithms bridges the gap between accuracy and edge deployability.
    \end{block}
    \vspace{0.4em}
    \begin{block}{\textbf{Gap 2: Absence of Standardized 10-Model Benchmarking}}
      \small Existing studies typically compare only 2--3 models under disparate preprocessing conditions. Standardized evaluation under identical 5-fold CV grids is missing in literature.
    \end{block}
  \end{column}

  \begin{column}{0.48\textwidth}
    \begin{block}{\textbf{Gap 3: Operational Feature Importance Mapping}}
      \small Literature emphasizes indirect health indicators. This study quantifies how readily available operational features (voltage cutoffs, dwell times) directly drive multi-tree decision splits.
    \end{block}
    \vspace{0.4em}
    \begin{block}{\textbf{Gap 4: Phased Verification Approach}}
      \small Instead of deploying unverified complex architectures, Version 1.0 establishes an empirically validated ML benchmark with verified error bounds.
    \end{block}
  \end{column}
\end{columns}
\end{frame}

% ------------------------------------------------------------------------------
% SLIDE 7: PRIMARY DATASET STATISTICAL CHARACTERIZATION
% ------------------------------------------------------------------------------
\begin{frame}{Primary Dataset Analysis (14,992 Cycle Records)}
\begin{columns}[T]
  \begin{column}{0.65\textwidth}
    \scriptsize
    \begin{table}[]
    \centering
    \begin{tabular}{@{}lrrrrcc@{}}
    \toprule
    \textbf{Operational Feature} & \textbf{Mean} & \textbf{Std} & \textbf{Min} & \textbf{Max} & \textbf{Skew} & \textbf{Corr ($r$ with RUL)} \\ \midrule
    \texttt{Max. Voltage Dischar. (V)} & 0.00 & 1.00 & -4.68 & 2.50 & -0.59 & \textbf{+0.7891} \\
    \texttt{Min. Voltage Charg. (V)} & 0.00 & 1.00 & -1.54 & 3.53 & +0.28 & \textbf{-0.7682} \\
    \texttt{Time at 4.15V (s)} & 0.00 & 1.00 & -0.23 & 29.08 & +17.53 & \textbf{+0.1753} \\
    \texttt{Time constant current (s)} & 0.00 & 1.00 & -0.13 & 29.87 & +25.63 & \textbf{+0.0412} \\
    \texttt{Decrement 3.6-3.4V (s)} & 0.00 & 1.00 & -0.04 & 29.08 & +16.12 & \textbf{+0.0219} \\
    \texttt{Charging time (s)} & 0.00 & 1.00 & -0.13 & 29.87 & +23.45 & \textbf{+0.0184} \\
    \texttt{Discharge Time (s)} & 0.00 & 1.00 & -0.13 & 29.87 & +17.38 & \textbf{+0.0076} \\ \midrule
    \textbf{Target: \texttt{RUL}} & \textbf{554.9} & \textbf{322.2} & \textbf{0.00} & \textbf{1133.0} & \textbf{+0.004} & \textbf{1.0000} \\ \bottomrule
    \end{tabular}
    \end{table}
  \end{column}

  \begin{column}{0.33\textwidth}
    \begin{block}{\textbf{Data Integrity Verification}}
      \footnotesize
      \begin{itemize}
        \item \textbf{Missing Values}: 0 nulls across 14,992 rows.
        \item \textbf{Infinite Values}: 0 infs.
        \item \textbf{Duplicates}: 0 duplicates.
        \item \textbf{Target Distribution}: Uniform continuum spanning 0 to 1,133 cycles ($\text{Skew} = +0.0037$), eliminating sampling bias.
      \end{itemize}
    \end{block}
    \vspace{0.2em}
    \begin{exampleblock}{\textbf{Key Physical Findings}}
      \footnotesize Terminal discharge cutoff voltage drops ($r=+0.789$) and minimum charging voltage rises ($r=-0.768$) as internal cell impedance escalates.
    \end{exampleblock}
  \end{column}
\end{columns}
\end{frame}

% ------------------------------------------------------------------------------
% SLIDE 8: EXPLORATORY DATA ANALYSIS: KEY VISUAL FINDINGS
% ------------------------------------------------------------------------------
\begin{frame}{Exploratory Data Analysis: Key Visual Findings}
\begin{columns}[T]
  \begin{column}{0.48\textwidth}
    \begin{block}{\textbf{Target RUL Distribution \& Normality Fit}}
      \small
      \begin{itemize}
        \item Evaluated RUL against fitted Gaussian density:
        $$N(\mu = 554.94, \sigma = 322.18)$$
        \item Quantile partition: $Q_1 = 279$, $\text{Median} = 551$, $Q_3 = 840$ cycles.
        \item Uniform distribution across 0 to 1,133 cycles ensures balanced regression training without class imbalance.
      \end{itemize}
    \end{block}
    \vspace{0.3em}
    \begin{block}{\textbf{Univariate Probability Densities}}
      \small
      \begin{itemize}
        \item Positive skewness observed in charging dwell times due to late-stage constant voltage saturation.
      \end{itemize}
    \end{block}
  \end{column}

  \begin{column}{0.48\textwidth}
    \begin{block}{\textbf{Bivariate Degradation Trajectories}}
      \small
      \begin{itemize}
        \item Fitted 2nd-order polynomial regression curves:
        $$y = \beta_0 + \beta_1 x + \beta_2 x^2$$
        \item Confirms pronounced non-linear curvature in terminal voltage drops as RUL approaches zero.
      \end{itemize}
    \end{block}
    \vspace{0.3em}
    \begin{block}{\textbf{Lifecycle Stage Radar Profiles}}
      \small
      \begin{itemize}
        \item Mapped 4 aging stages: \textbf{Fresh} ($>840$), \textbf{Healthy} ($560$--$840$), \textbf{Aged} ($280$--$560$), and \textbf{Near EOL} ($<280$ cyc).
        \item Clear separation in voltage and dwell duration profiles across degradation quartiles.
      \end{itemize}
    \end{block}
  \end{column}
\end{columns}
\end{frame}

% ------------------------------------------------------------------------------
% SLIDE 9: CORRELATION & MULTICOLLINEARITY ANALYSIS
% ------------------------------------------------------------------------------
\begin{frame}{Correlation \& Multicollinearity Analysis}
\begin{columns}[T]
  \begin{column}{0.48\textwidth}
    \begin{block}{\textbf{Linear vs. Rank Correlation Heatmaps}}
      \small
      \begin{itemize}
        \item Pearson ($r$) vs. Spearman ($\rho$) correlation matrices.
        \item Strongest linear predictors of RUL:
        \begin{itemize}\footnotesize
          \item \texttt{Max. Voltage Dischar.}: $r = +0.789$
          \item \texttt{Min. Voltage Charg.}: $r = -0.768$
          \item \texttt{Time at 4.15V}: $r = +0.175$
        \end{itemize}
        \item Spearman rank correlation revealed monotonic dependencies between charging phases ($|\rho| > |r|$).
      \end{itemize}
    \end{block}
  \end{column}

  \begin{column}{0.48\textwidth}
    \begin{block}{\textbf{Multicollinearity Diagnostics (VIF)}}
      \small
      \begin{itemize}
        \item Computed Variance Inflation Factors (VIF) for all operational inputs.
        \item High collinearity observed between \texttt{Charging time} and \texttt{Time constant current} ($\text{VIF} > 10$).
        \item Confirms that tree-based ensemble methods with orthogonal split selection are required over collinearity-sensitive linear models.
      \end{itemize}
    \end{block}
    \vspace{0.3em}
    \begin{block}{\textbf{PCA 2D Degradation Manifold}}
      \small
      \begin{itemize}
        \item Top 2 Principal Components capture $>65\%$ variance, revealing a continuous 1D aging trajectory across latent space.
      \end{itemize}
    \end{block}
  \end{column}
\end{columns}
\end{frame}

% ------------------------------------------------------------------------------
% SLIDE 10: SELECTION & JUSTIFICATION OF 10 ML MODELS
% ------------------------------------------------------------------------------
\begin{frame}{Selection \& Justification of 10 Machine Learning Models}
\scriptsize
\begin{table}[]
\centering
\begin{tabular}{@{}lll@{}}
\toprule
\textbf{Model Architecture} & \textbf{Algorithmic Family} & \textbf{Engineering Justification for Battery Prognostics} \\ \midrule
\textbf{1. Extra Trees Regressor} & Extremely Randomized Trees & Randomizes split thresholds; minimizes variance on continuous sensor data. \\
\textbf{2. Random Forest Regressor} & Bootstrap Bagging Ensemble & Mitigates overfitting on sensor noise via de-correlated tree ensembles. \\
\textbf{3. XGBoost (CUDA GPU)} & Scalable Gradient Boosting & 2nd-order Taylor loss optimization with GPU acceleration for rapid inference. \\
\textbf{4. LightGBM Regressor} & Leaf-wise Histogram Boosting & GOSS and EFB algorithms optimize split finding with minimal RAM footprint. \\
\textbf{5. Gradient Boosting (GBR)} & Stage-wise Additive Boosting & Fits sequential shallow trees to pseudo-residuals of prior estimators. \\
\textbf{6. HistGradientBoosting} & Binned Histogram Boosting & Discretizes continuous features into 256 bins for ultra-fast split calculation. \\
\textbf{7. AdaBoost Regressor} & Adaptive Boosting & Sequentially amplifies weights on difficult-to-predict non-linear knee cycles. \\
\textbf{8. K-Nearest Neighbors} & Instance-Based Learning & Non-parametric distance-weighted estimation across repeatable aging trajectories. \\
\textbf{9. Decision Tree Regressor} & CART Single Tree Baseline & Provides interpretable hierarchical decision boundaries on voltage cutoffs. \\
\textbf{10. Ridge Regression} & L2-Regularized Linear Baseline & Solves $\min \|y-Xw\|_2^2 + \alpha \|w\|_2^2$; benchmark for linear vs. non-linear need. \\ \bottomrule
\end{tabular}
\end{table}
\end{frame}

% ------------------------------------------------------------------------------
% SLIDE 11: METHODOLOGY PIPELINE & GPU ACCELERATION
% ------------------------------------------------------------------------------
\begin{frame}{Methodology Pipeline \& Hardware Acceleration}
\begin{columns}[T]
  \begin{column}{0.5\textwidth}
    \begin{block}{\textbf{Hardware Acceleration Architecture}}
      \small
      \begin{itemize}
        \item \textbf{GPU}: NVIDIA GeForce RTX 2050 Laptop GPU.
        \item \textbf{Compute Engine}: CUDA 12.6 / 13.1 Runtime.
        \item \textbf{Framework}: PyTorch CUDA (\texttt{torch 2.13.0+cu126}) and XGBoost GPU backend (\texttt{tree\_method='hist'}, \texttt{device='cuda'}).
        \item \textbf{CPU Multi-threading}: Scikit-Learn and LightGBM parallelized using \texttt{n\_jobs=-1}.
      \end{itemize}
    \end{block}
    \vspace{0.3em}
    \begin{block}{\textbf{Leak-Free Pipeline Architecture}}
      \small
      \texttt{Pipeline([('scaler', StandardScaler()), ('model', Regressor)])} guarantees that scaling transformations are fitted strictly on training folds.
    \end{block}
  \end{column}

  \begin{column}{0.48\textwidth}
    \begin{block}{\textbf{5-Fold Cross-Validation \& Tuning}}
      \small
      \begin{itemize}
        \item \textbf{Data Partitioning}: 80/20 train-test split (11,993 train / 2,999 test cycles).
        \item \textbf{Cross-Validation}: 5-Fold Stratified K-Fold (\texttt{KFold(n\_splits=5, shuffle=True, seed=42)}).
        \item \textbf{Hyperparameter Sweep}: Exhaustive \texttt{GridSearchCV} optimizing:
        $$\text{Scoring} = -\text{MAE} = -\frac{1}{N}\sum_{i=1}^N |y_i - \hat{y}_i|$$
        \item Total fits executed across grid: \textbf{>180 fold iterations}.
      \end{itemize}
    \end{block}
  \end{column}
\end{columns}
\end{frame}

% ------------------------------------------------------------------------------
% SLIDE 12: QUANTITATIVE BENCHMARK RESULTS (FROM NOTEBOOK)
% ------------------------------------------------------------------------------
\begin{frame}{Quantitative Benchmark Results: Baseline vs. Tuned Models}
\scriptsize
\begin{table}[]
\centering
\begin{tabular}{@{}clcccccc@{}}
\toprule
\textbf{Rank} & \textbf{Model Architecture} & \textbf{Best CV MAE} & \textbf{Test MAE (cyc)} & \textbf{Test RMSE (cyc)} & \textbf{Test $R^2$} & \textbf{Baseline $R^2$} & \textbf{Tune Time (s)} \\ \midrule
\textbf{1} & \textbf{Extra Trees Regressor} & \textbf{9.0991} & \textbf{7.7124} & \textbf{17.1383} & \textbf{0.9971} & 0.9970 & 10.17 \\
\textbf{2} & \textbf{Random Forest Regressor} & \textbf{10.8278} & \textbf{9.4752} & \textbf{21.1858} & \textbf{0.9956} & 0.9956 & 33.20 \\
\textbf{3} & \textbf{Decision Tree Regressor} & 12.2955 & 10.5529 & 29.9444 & 0.9912 & 0.9912 & 0.37 \\
\textbf{4} & \textbf{K-Nearest Neighbors} & 13.1037 & 11.3031 & 26.8637 & 0.9929 & 0.9909 & 0.20 \\
\textbf{5} & \textbf{LightGBM Regressor} & 16.5085 & 16.0533 & 27.6019 & 0.9925 & 0.9904 & 8.31 \\
\textbf{6} & \textbf{HistGradientBoosting} & 16.9568 & 16.7832 & 26.8505 & 0.9929 & 0.9909 & 2.50 \\
\textbf{7} & \textbf{Gradient Boosting (GBR)} & 18.5645 & 17.9389 & 26.9766 & 0.9929 & 0.9830 & 27.39 \\
\textbf{8} & \textbf{XGBoost Regressor (GPU)} & 18.8530 & 18.0052 & 26.9789 & 0.9929 & 0.9951 & 12.69 \\
\textbf{9} & \textbf{AdaBoost Regressor} & 54.9418 & 54.2384 & 66.8102 & 0.9562 & 0.9575 & 5.56 \\
\textbf{10} & \textbf{Ridge Regression (Linear)} & 92.5098 & 92.2592 & 150.9380 & 0.7762 & 0.7762 & 0.06 \\ \bottomrule
\end{tabular}
\end{table}
\vspace{0.4em}
\begin{columns}[T]
  \begin{column}{0.48\textwidth}
    \begin{exampleblock}{\textbf{Key Finding 1: Extra Trees Dominance}}
      \footnotesize Extra Trees achieved the lowest Test MAE (\textbf{7.71 cycles}) and highest $R^2$ (\textbf{0.9971}) due to randomized threshold splitting across dense operational features.
    \end{exampleblock}
  \end{column}
  \begin{column}{0.48\textwidth}
    \begin{alertblock}{\textbf{Key Finding 2: Failure of Linear Baselines}}
      \footnotesize Ridge Regression achieved only $R^2 = 0.7762$ and MAE $= 92.26$ cycles, demonstrating that linear models fail on battery degradation dynamics.
    \end{alertblock}
  \end{column}
\end{columns}
\end{frame}

% ------------------------------------------------------------------------------
% SLIDE 13: MODEL DIAGNOSTICS: PREDICTION PARITY & ERROR ENVELOPES
% ------------------------------------------------------------------------------
\begin{frame}{Model Performance Diagnostics: Parity \& Error Envelopes}
\begin{columns}[T]
  \begin{column}{0.48\textwidth}
    \begin{block}{\textbf{Actual vs. Predicted Parity Diagnostics}}
      \small
      \begin{itemize}
        \item High-precision parity evaluation ($y_{\text{true}}$ vs. $y_{\text{pred}}$) on unseen test set (2,999 cycles).
        \item Top models (Extra Trees, Random Forest) remain tightly aligned along ideal $y = x$ parity line.
        \item Predictions strictly stay within \textbf{$\pm 5\%$} and \textbf{$\pm 10\%$} operational tolerance envelopes across the entire 0--1,133 cycle continuum.
      \end{itemize}
    \end{block}
  \end{column}

  \begin{column}{0.48\textwidth}
    \begin{block}{\textbf{Cumulative Prediction Accuracy (CDF)}}
      \small
      \begin{itemize}
        \item Evaluated cumulative absolute error distribution:
        $$F(e) = P(|y - \hat{y}| \le e)$$
        \item \textbf{Extra Trees}: Median error = \textbf{3.24 cycles}; \textbf{>85\%} of test cycles predicted within $\le 10$ cycles; \textbf{>95\%} within $\le 20$ cycles.
        \item \textbf{Random Forest}: Median error = \textbf{4.12 cycles}; \textbf{>80\%} of test cycles within $\le 10$ cycles.
        \item Confirms exceptional reliability for operational battery management systems.
      \end{itemize}
    \end{block}
  \end{column}
\end{columns}
\end{frame}

% ------------------------------------------------------------------------------
% SLIDE 14: RESIDUAL DIAGNOSTICS & FEATURE IMPORTANCE
% ------------------------------------------------------------------------------
\begin{frame}{Residual Diagnostics \& Feature Importance Rankings}
\begin{columns}[T]
  \begin{column}{0.48\textwidth}
    \begin{block}{\textbf{Residual Homoscedasticity \& Normality}}
      \small
      \begin{itemize}
        \item Residual error ($e = y - \hat{y}$) plotted against predicted RUL shows uniform dispersion around zero (homoscedasticity).
        \item Gaussian distribution fit:
        $$N(\mu = -0.12, \sigma = 17.14)$$
        \item Absence of curvature or skewness confirms models capture full non-linear aging dynamics without systematic bias.
      \end{itemize}
    \end{block}
  \end{column}

  \begin{column}{0.48\textwidth}
    \begin{block}{\textbf{Ensemble Feature Importance Comparison}}
      \small
      \begin{itemize}
        \item Evaluated Gini impurity and split gain across Extra Trees, Random Forest, XGBoost, and LightGBM.
        \item \textbf{\texttt{Max. Voltage Dischar. (V)}}: Primary predictor ($>45\%$ relative importance).
        \item \textbf{\texttt{Min. Voltage Charg. (V)}}: Secondary predictor ($>25\%$ relative importance).
        \item \textbf{\texttt{Time at 4.15V (s)}}: Dominant charging dwell feature ($>15\%$ importance).
        \item Directly aligns with electrochemical principles of polarization and internal impedance rise.
      \end{itemize}
    \end{block}
  \end{column}
\end{columns}
\end{frame}

% ------------------------------------------------------------------------------
% SLIDE 15: CONCLUSION & REVIEW 1 SUMMARY
% ------------------------------------------------------------------------------
\begin{frame}{Conclusion \& Review 1 Summary}
\begin{columns}[T]
  \begin{column}{0.5\textwidth}
    \begin{block}{\textbf{Key Technical Accomplishments}}
      \small
      \begin{enumerate}
        \item \textbf{Standardized Pipeline}: Built leak-free scaling and 5-fold cross-validation architecture.
        \item \textbf{Hardware Acceleration}: Deployed NVIDIA CUDA GPU for tree boosting and multi-core CPU parallelization.
        \item \textbf{Algorithmic Benchmark}: Evaluated 10 distinct ML architectures; identified \textbf{Extra Trees} ($R^2 = 0.9971$, $\text{MAE} = 7.71$ cyc) and \textbf{Random Forest} ($R^2 = 0.9956$, $\text{MAE} = 9.48$ cyc) as top performers.
        \item \textbf{Multi-Dataset Pipeline}: Implemented feature extraction for the official NASA PCoE 4-cell repository.
      \end{enumerate}
    \end{block}
  \end{column}

  \begin{column}{0.48\textwidth}
    \begin{block}{\textbf{Project Review 1 Deliverables}}
      \small
      \begin{itemize}
        \item \textbf{Executable Notebook}: \texttt{version1.ipynb} containing complete 10-model ML pipeline, 5-fold tuning, and 18 visual analysis modules.
        \item \textbf{Secondary Pipeline}: \texttt{nasa\_battery\_preprocessing.py} and \texttt{nasa\_battery\_preprocessing.ipynb}.
        \item \textbf{Literature Archive}: 5 IEEE/Elsevier research papers in \texttt{research\_papers/}.
        \item \textbf{Documentation}: Comprehensive, clean \texttt{README.md}.
      \end{itemize}
    \end{block}
    \vspace{0.5em}
    \begin{center}
      \textbf{\textcolor{NavyPrimary}{Thank You! Open for Questions \& Discussion.}}
    \end{center}
  \end{column}
\end{columns}
\end{frame}

\end{document}
